\begin{proof}
We use the exact scalar star theorem
\Cref{thm:p071-infinite-star-collapse-candidate} and the compact-core
envelope lemmas
\Cref{lem:p071-exact-star-polyhedral-envelope,%
lem:p071-arbitrary-exact-star-polyhedral-insertion}.

\medskip\noindent
\emph{The comparison sheaf and the easy inclusion.}
Upper semicontinuity, branchwise joint convexity, and the endpoint trace are
local.  The junction inequality is local by the same contraction calculation
as in the open-Euclidean product.  Namely, contract every sample by a common
factor \(0<s<1\) towards \((\eta,y_0)\):
\[
 (-r_i,y_i)\longmapsto
       (-sr_i,\,y_0+s(y_i-y_0)).
\]
The barycentre relation is preserved and every reciprocal weight is
multiplied by \(s^{-1}\).  If restriction makes finitely many formerly full
branches exceptional, add samples on them with polyhedral coordinate
\(y_0\).  Their values are at most the centre value.  Branchwise convexity
bounds every contracted original sample by
\[
 (1-s)c(y_0)+s\,u_i(-r_i,y_i).
\]
Substitution in the restricted-chart inequality and cancellation of the
added central terms recover
\eqref{eq:p071-infinite-polyhedral-junction} on the original chart.  Thus
\(\mathcal J_X\) is a sheaf.

It is closed under pointwise decreasing sequential limits.  Such a limit is
upper semicontinuous and jointly extended-real convex on every branch.
The exact full-endpoint trace passes by monotonicity: one inequality is
immediate, while one sufficiently advanced row gives the reverse inequality,
also when the endpoint value is \(-\infty\).  At a node, fix a compact
polytope \(Q\) in one comparison chart.  For each row, finitely many charts
cover \(\{\eta\}\times Q\); the union of their finite exceptional and
countable activity sets is countable.  Taking the union over the rows gives
one countable activity set for the limit.

If the limiting centre value in a prescribed junction tuple is finite, then
the centre values of all rows are finite.  The rowwise junction inequality
forces every sampled row value to be finite, since no term may equal
\(+\infty\).  None of these sampled values can tend to \(-\infty\): all
other terms are bounded above by their first-row values, so the right side
would tend to \(-\infty\) while the left side has a finite limit.  Termwise
passage to the limit proves the junction inequality.  Hence
\(\mathcal J_X\) is decreasing-sequentially closed.

Every generator belongs to \(\mathcal J_X\).  Only the junction inequality
needs comment.  Put \(d_i=\partial_tg_i(0,y_0)\).  Joint convexity gives
\[
 g_i(-r_i,y_i)\geq
 c_g(y_0)-r_id_i+
 \nabla c_g(y_0)\mathbin\cdot(y_i-y_0).
\]
After division by \(r_i\), the polyhedral linear terms cancel by
\eqref{eq:p071-infinite-polyhedral-barycentre}.  Every omitted full-branch
derivative is nonnegative, and
\eqref{eq:p071-infinite-polyhedral-smooth-balance} therefore makes the sum
of the selected derivatives nonpositive.  This is exactly
\eqref{eq:p071-infinite-polyhedral-junction}.  Minimality gives
\begin{equation}
                           \mathcal P_X\subseteq\mathcal J_X.
 \label{eq:p071-infinite-polyhedral-easy-inclusion}
\end{equation}

\medskip\noindent
\emph{Continuous decreasing envelope rows at a node.}
Fix \(u\in\mathcal J_X\) near \((\eta,y_0)\).  Choose a comparison chart
\(U\times V_0\), with finite exceptional set \(A\) and countable activity
set \(K\).  Shrink the exceptional collars once and let \(D\subset U\) be
the resulting compact closed standard core.  Choose a relatively open
convex set \(V\subset P\) and a compact convex polytope \(Q\subset P\) such
that
\[
            y_0\in V,\qquad
            \overline V^{\,P}\subset Q\subset V_0.
\]
Collapse every full branch outside \(A\cup K\) to the node.  The resulting
compact exact star \(D_K\) is metrizable, and \(u|_{D\times Q}\) factors
through \(D_K\times Q\).

Fix a compatible metric \(d\) on this compact product and set
\begin{equation}
 \phi_n(z)=\sup_w
   \bigl(\max\{u(w),-n\}-n\,d(z,w)\bigr).
 \label{eq:p071-infinite-polyhedral-upper-regularization}
\end{equation}
These functions are finite and continuous,
\(\phi_{n+1}\leq\phi_n\), and \(\phi_n\downarrow u\).  Define
\[
                         H_n=C^{\mathrm{ex}}\phi_n.
\]
The compact-envelope lemma gives a finite continuous exact comparison
function and monotonicity in the obstacle.  The insertion lemma gives
\(u\leq H_n\), while the defining minorant property gives
\(H_n\leq\phi_n\).  Therefore
\begin{equation}
             H_{n+1}\leq H_n,\qquad H_n\downarrow u
             \quad\text{on }D_K\times Q.
 \label{eq:p071-infinite-polyhedral-envelope-rows}
\end{equation}
Lift \(H_n\) to \(D\times Q\) by making every branch outside \(K\) equal
to its node slice.  The two-stage barycentre calculation in
\Cref{lem:p071-arbitrary-exact-star-polyhedral-insertion} shows that the
lift is still an exact product comparison function.  Equations
\eqref{eq:p071-infinite-polyhedral-envelope-rows} remain valid after the
lift.

\medskip\noindent
\emph{Finite activity and a uniform smooth majorant.}
We first record a uniform approximation statement for one finite continuous
row \(H\) on \(D\times Q\).  Since a full branch is nonincreasing outwards,
\[
                         H_b(t,y)\leq H(\eta,y).
\]
Continuity at the node in the cofinite-branch topology, uniformly over the
compact set \(Q\), implies that for every \(\varepsilon>0\) there is a finite
set \(F_\varepsilon\subset K\) such that replacing every branch outside
\(A\cup F_\varepsilon\) by the node slice produces a finite-active comparison
function \(\widetilde H\) satisfying
\begin{equation}
                 H\leq\widetilde H\leq H+\varepsilon
                 \quad\text{on }D\times Q.
 \label{eq:p071-infinite-polyhedral-finite-activity-error}
\end{equation}
The junction condition survives this replacement even when an inactive
sample has a different polyhedral coordinate: first apply the finite-substar
junction inequality at the barycentre of the exceptional and active
samples, then combine that point with the inactive coordinates by convexity
of the node slice.

Choose once a smaller standard neighbourhood \(U'\) whose closure is
contained in the interior of the chosen core.  We claim that there is an
ambient function \(W\) on \(U'\times L_P\), whose restriction to
\(U'\times V\) is a generator, such that
\begin{equation}
               \widetilde H\leq W\leq\widetilde H+5\varepsilon
               \quad\text{on }\overline{U'}\times Q.
 \label{eq:p071-infinite-polyhedral-smooth-majorant}
\end{equation}
Choose an origin in \(L_P\).  For \(p\in L_P^*\), put
\begin{align}
 f_0(p)&=\min_{y\in Q}
       \bigl(\widetilde H(\eta,y)-p\mathbin\cdot y\bigr),\\
 f_b(t,p)&=\min_{y\in Q}
       \bigl(\widetilde H_b(t,y)-p\mathbin\cdot y\bigr).
 \label{eq:p071-infinite-polyhedral-scalar-profiles}
\end{align}
For each \(p\), these functions form a finite, continuous, finite-active
member of the scalar exact-star comparison class.  Branch convexity follows
by combining minimizing polyhedral points.  The endpoint trace passes
through the minimum by uniform continuity on the compact terminal prism.
For the scalar junction inequality, choose minimizing points for the finitely
many branch terms and use their reciprocal-weight barycentre as a competitor
in the node minimum.

For fixed \(b\) and \(t\), extend the slice
\(\widetilde H_b(t,\cdot)\) by \(+\infty\) off \(Q\).  The extension is
proper, lower semicontinuous, and convex on \(L_P\), including at every face
and corner because \(\widetilde H\) is continuous on the closed polytope.
Fenchel--Moreau gives
\begin{equation}
 \widetilde H_b(t,y)=
       \sup_{p\in L_P^*}\bigl(p\mathbin\cdot y+f_b(t,p)\bigr)
 \qquad(y\in Q),
 \label{eq:p071-infinite-polyhedral-fenchel}
\end{equation}
with the analogous formula at the node and full endpoints.

At each point of the compact product
\(\overline{U'}\times Q\), choose a covector realizing
\eqref{eq:p071-infinite-polyhedral-fenchel} within \(\varepsilon\).
For fixed \(p\), the profile \(f(\,\cdot\,,p)\) is continuous on the star,
so one covector works on a neighbourhood.  A finite subcover gives
\(p_1,\ldots,p_q\) such that
\begin{equation}
 \widetilde H-\varepsilon\leq
 A_0:=\max_j\bigl(p_j\mathbin\cdot y+f(x,p_j)\bigr)
 \leq\widetilde H.
 \label{eq:p071-infinite-polyhedral-finite-dual-sampling}
\end{equation}

For each selected covector, the finite harmonic smoothing in the scalar
exact-star theorem gives a smooth psh profile \(s_j\) on the same \(U'\)
with
\[
             \lVert s_j-f(\,\cdot\,,p_j)\rVert_{\overline{U'}}
                    <\varepsilon.
\]
Lift these profiles to
\[
               s_j(x)+p_j\mathbin\cdot y+2\varepsilon.
\]
Their maximum lies between \(\widetilde H\) and
\(\widetilde H+3\varepsilon\).  Apply a smooth convex coordinatewise
nondecreasing regularized maximum with error at most \(\varepsilon\), and
add one further \(\varepsilon\).  The result is \(W\) and satisfies
\eqref{eq:p071-infinite-polyhedral-smooth-majorant}.  It is an exact product
generator.  Joint convexity follows from convex monotone composition.  At
the node, the derivatives of the regularized maximum with respect to its
inputs are branch-independent, so the sum over branches reduces to the
exact balance of each \(s_j\).  Only finitely many branches are active, and
the finitely many scalar profiles have one common terminal collar on each
full active branch.

Apply this construction to each \(H_n\).  Fix a summable positive sequence
\((\varepsilon_n)\), choose \(\widetilde H_n\) as in
\eqref{eq:p071-infinite-polyhedral-finite-activity-error}, and then choose
\(W_n\) as in
\eqref{eq:p071-infinite-polyhedral-smooth-majorant}.  Thus
\begin{equation}
                    H_n\leq W_n\leq H_n+6\varepsilon_n
                    \quad\text{on }\overline{U'}\times Q.
 \label{eq:p071-infinite-polyhedral-row-error}
\end{equation}
Put
\[
             d_n=6\sum_{k>n}\varepsilon_k,\qquad G_n=W_n+d_n.
\]
Using \eqref{eq:p071-infinite-polyhedral-envelope-rows} and
\eqref{eq:p071-infinite-polyhedral-row-error}, we obtain
\[
 G_{n+1}\leq H_{n+1}+6\varepsilon_{n+1}+d_{n+1}
       \leq H_n+d_n\leq W_n+d_n=G_n.
\]
Moreover \(G_n\to u\) pointwise.  This proves the first-stage assertion on
one fixed node neighbourhood, without a separate finite-centre case.

\medskip\noindent
\emph{Full endpoints and ordinary branch points.}
At a full endpoint, work on a compact terminal product
\([-\infty,R]\times Q\), with \(R<0\), inside the given chart.  The metric
upper regularizations
\eqref{eq:p071-infinite-polyhedral-upper-regularization}, followed by the
terminal insertion operator of
\Cref{lem:p071-terminal-polyhedral-insertion}, give finite continuous
jointly convex rows \(H_n\downarrow u\), nondecreasing in the branch
parameter and with exact endpoint trace.

For one such row, take the same partial Legendre profiles
\[
                 f(t,p)=\min_{y\in Q}(H(t,y)-p\mathbin\cdot y).
\]
Continuity on the closed polytope gives boundary-valid Fenchel--Moreau
reconstruction and finite dual sampling on a fixed smaller terminal core.
For each selected covector, the scalar endpoint theorem gives a smooth
convex nondecreasing profile, exactly constant on a terminal collar, within
\(\varepsilon\) of \(f(\,\cdot\,,p)\).  The lifted regularized-maximum
construction therefore gives \(W_n\) with
\[
                         H_n\leq W_n\leq H_n+5\varepsilon_n
\]
on one fixed terminal product neighbourhood.  The same summable-tail
calculation, with coefficient \(5\), makes the resulting
generators decrease to \(u\).

At an ordinary finite point of a branch there is no junction.  A smaller
chart is a relatively open convex subset of the polyhedron
\(\mathbf R\times P\), and the one-point case of the polyhedral convex
collapse gives a decreasing sequence of ambient smooth jointly convex
functions on one fixed smaller chart.

We have proved
\(\mathcal J_X\subseteq\mathcal P_{X,1}\) at every point.  Together with
\eqref{eq:p071-infinite-polyhedral-easy-inclusion} and the tautological
inclusion
\(\mathcal P_{X,1}\subseteq\mathcal P_X\), this proves
\eqref{eq:p071-infinite-polyhedral-collapse} and all fixed-neighbourhood
assertions.
\end{proof}
